\reftitle{References}
\begin{thebibliography}{99}
\bibitem{ref-cdvae} Xie T, Fu X, Ganea OE, Barzilay R, Jaakkola T. Crystal diffusion variational autoencoder for periodic material generation. Paper presented at: 10th International Conference on Learning Representations; 2022 Apr 25--29; Virtual.
\bibitem{ref-diffcsp} Jiao R, Huang W, Lin P, Han J, Chen P, Lu Y, et al. Crystal structure prediction by joint equivariant diffusion. Adv Neural Inf Process Syst. 2023;36:17464--97.
\bibitem{ref-mattergen} Zeni C, Pinsler R, Zuegner D, Fowler A, Horton M, Fu X, et al. A generative model for inorganic materials design. Nature. 2025;639:624--32. doi: 10.1038/s41586-025-08628-5.
\bibitem{ref-chemeleon} Park H, Onwuli A, Walsh A. Exploration of crystal chemical space using text-guided generative artificial intelligence. Nat Commun. 2025;16(1):4379. doi: 10.1038/s41467-025-59636-y.
\bibitem{ref-materials-project} Jain A, Ong SP, Hautier G, Chen W, Richards WD, Dacek S, et al. Commentary: The Materials Project: a materials genome approach to accelerating materials innovation. APL Mater. 2013;1(1):011002. doi: 10.1063/1.4812323.
\bibitem{ref-oqmd} Kirklin S, Saal JE, Meredig B, Thompson A, Doak JW, Aykol M, et al. The Open Quantum Materials Database (OQMD): assessing the accuracy of DFT formation energies. npj Comput Mater. 2015;1:15010. doi: 10.1038/npjcompumats.2015.10.
\bibitem{ref-aflow} Curtarolo S, Setyawan W, Hart GLW, Jahnatek M, Chepulskii RV, Taylor RH, et al. AFLOW: an automatic framework for high-throughput materials discovery. Comput Mater Sci. 2012;58:218--26. doi: 10.1016/j.commatsci.2012.02.005.
\bibitem{ref-jarvis} Choudhary K, Garrity KF, Reid ACE, DeCost B, Biacchi AJ, Hight Walker AR, et al. The joint automated repository for various integrated simulations (JARVIS) for data-driven materials design. npj Comput Mater. 2020;6:173. doi: 10.1038/s41524-020-00440-1.
\bibitem{ref-badge} Ash JT, Zhang C, Krishnamurthy A, Langford J, Agarwal A. Deep batch active learning by diverse, uncertain gradient lower bounds. Paper presented at: 8th International Conference on Learning Representations; 2020 Apr 26--30; Virtual.
\bibitem{ref-coreset} Sener O, Savarese S. Active learning for convolutional neural networks: a core-set approach. Paper presented at: 6th International Conference on Learning Representations; 2018 Apr 30--May 3; Vancouver, Canada.
\bibitem{ref-kmeanspp} Arthur D, Vassilvitskii S. k-means++: the advantages of careful seeding. In: Proceedings of the 18th Annual ACM-SIAM Symposium on Discrete Algorithms; 2007 Jan 7--9; New Orleans, LA, USA. p. 1027--35.
\bibitem{ref-localpenalization} Gonzalez J, Dai Z, Hennig P, Lawrence N. Batch Bayesian optimization via local penalization. In: Proceedings of the 19th International Conference on Artificial Intelligence and Statistics; 2016 May 9--11; Cadiz, Spain. Proc Mach Learn Res. 2016;51:648--57.
\bibitem{ref-cgcnn} Xie T, Grossman JC. Crystal graph convolutional neural networks for an accurate and interpretable prediction of material properties. Phys Rev Lett. 2018;120(14):145301. doi: 10.1103/PhysRevLett.120.145301.
\bibitem{ref-alignn} Choudhary K, DeCost B. Atomistic line graph neural network for improved materials property predictions. npj Comput Mater. 2021;7:185. doi: 10.1038/s41524-021-00650-1.
\bibitem{ref-chgnet} Deng B, Zhong P, Jun K, Riebesell J, Han K, Bartel CJ, et al. CHGNet as a pretrained universal neural network potential for charge-informed atomistic modelling. Nat Mach Intell. 2023;5:1031--41. doi: 10.1038/s42256-023-00716-3.
\bibitem{ref-macemp} Batatia I, Benner P, Chiang Y, Elena AM, Kovacs DP, Riebesell J, et al. A foundation model for atomistic materials chemistry. J Chem Phys. 2025;163:184110. doi: 10.1063/5.0297006.
\bibitem{ref-mcdropout} Gal Y, Ghahramani Z. Dropout as a Bayesian approximation: representing model uncertainty in deep learning. In: Proceedings of the 33rd International Conference on Machine Learning; 2016 Jun 20--22; New York, NY, USA. Proc Mach Learn Res. 2016;48:1050--59.
\bibitem{ref-vasp} Kresse G, Furthm\"uller J. Efficient iterative schemes for ab initio total-energy calculations using a plane-wave basis set. Phys Rev B. 1996;54(16):11169--86. doi: 10.1103/PhysRevB.54.11169.
\bibitem{ref-paw} Bl\"ochl PE. Projector augmented-wave method. Phys Rev B. 1994;50(24):17953--79. doi: 10.1103/PhysRevB.50.17953.
\bibitem{ref-paw-vasp} Kresse G, Joubert D. From ultrasoft pseudopotentials to the projector augmented-wave method. Phys Rev B. 1999;59(3):1758--75. doi: 10.1103/PhysRevB.59.1758.
\bibitem{ref-pbe} Perdew JP, Burke K, Ernzerhof M. Generalized gradient approximation made simple. Phys Rev Lett. 1996;77(18):3865--68. doi: 10.1103/PhysRevLett.77.3865.
\bibitem{ref-dudarev} Dudarev SL, Botton GA, Savrasov SY, Humphreys CJ, Sutton AP. Electron-energy-loss spectra and the structural stability of nickel oxide: an LSDA+$U$ study. Phys Rev B. 1998;57(3):1505--09. doi: 10.1103/PhysRevB.57.1505.
\bibitem{ref-randomforest} Breiman L. Random forests. Mach Learn. 2001;45(1):5--32. doi: 10.1023/A:1010933404324.
\bibitem{ref-gradientboosting} Friedman JH. Greedy function approximation: a gradient boosting machine. Ann Stat. 2001;29(5):1189--1232. doi: 10.1214/aos/1013203451.
\bibitem{ref-sklearn} Pedregosa F, Varoquaux G, Gramfort A, Michel V, Thirion B, Grisel O, et al. Scikit-learn: machine learning in Python. J Mach Learn Res. 2011;12:2825--30.
\end{thebibliography}
